\section{Discussione: Sintesi Trasversale, Implicazioni Manageriali e Limiti}
\label{sec:discussion}

\subsection{Sintesi Trasversale dei Risultati}
La lettura integrata delle viste sviluppate in DashMe consente di superare l'analisi per singoli indicatori isolati, evidenziando tre dinamiche strutturali tra loro interconnesse. In primo luogo, emerge una chiara dissociazione tra \textit{volume assoluto} e \textit{criticità clinica}: i comparti dei dispositivi medici e dei prodotti farmaceutici guidano i volumi complessivi di richiamo presso grandi operatori industriali, ma rimangono ancorati prevalentemente alla Classe~II (rispettivamente 88,33\% e 80,76\%); di converso, i comparti \textit{Prodotti per Animali} (47,65\% Classe~I) e \textit{Cibo e Bevande} ($\approx$ 43,57\% Classe~I) concentrano la quota proporzionalmente più severa di eventi potenzialmente letali.

In secondo luogo, la dimensione geografica e quella della portata distributiva mostrano come il rischio non resti confinato all'area di origine. Se gli Stati di California, Texas, Florida, il Nord-Est/Midwest e Puerto Rico rappresentano i principali nodi di notifica (fino a 4.522 eventi per singolo Stato), la matrice di Dashboard~3 conferma che una quota rilevante di richiami di Classe~I raggiunge circuiti \textit{Regionali (Multi-Stato)} e \textit{Internazionali / Export}, amplificando la complessità del coordinamento inter-giurisdizionale.

In terzo luogo, l'osservazione congiunta delle cause di richiamo, della conservazione e del confezionamento suggerisce che contaminazioni microbiologiche (\textit{Salmonella}, \textit{Listeria}, \textit{E.Coli}) e allergeni non dichiarati trovano un vettore ricorrente negli imballaggi in \textit{Plastica / Busta} (oltre 8.000 prodotti) e \textit{Vetro / Bottiglia / Barattolo} (circa 4.800 prodotti), interessando sia il segmento a temperatura ambiente sia le filiere \textit{Refrigerato / Fresco} e \textit{Surgelato}. \todo{Integrare la sintesi trasversale con le dinamiche temporali e la classifica degli allergeni non appena completata l'analisi della Dashboard 2.}

\subsection{Implicazioni per il Project Management e il Risk Management}
Dal punto di vista gestionale e dell'ingegneria dei processi di progetto, i risultati offrono indicazioni operative concrete per i responsabili della qualità, della \textit{supply chain} e della conformità regolatoria:
\begin{enumerate}
    \item \textbf{Audit selettivi sui fornitori e allocazione delle risorse di controllo}: la matrice di rischio emersa dalla Dashboard~1 suggerisce di differenziare la frequenza e la profondità degli audit in base al profilo di letalità settoriale e alla concentrazione territoriale. Per fornitori attivi in \textit{Cibo e Bevande}, \textit{Prodotti per Animali} o presso siti ad alta incidenza di Classe~I, gli audit dovrebbero privilegiare test microbiologici e controlli sui cambi formato/etichetta; nei dispositivi medici e nei centri ematologici, invece, l'attenzione va posta sul controllo statistico di processo e sulla gestione documentale delle non conformità ricorrenti.
    \item \textbf{Governance della logistica inversa e chiusura dei 18.555 richiami ``In corso''}: la presenza di 18.555 prodotti in stato ``In corso'' su 104.537 totali rilevati nella Dashboard~3 evidenzia come l'apertura formale di un richiamo non garantisca una rapida chiusura operativa. Nei progetti di gestione delle crisi (\textit{recall management}), è indispensabile trattare ogni evento come un programma strutturato con \textit{milestone} di tracciabilità, riconciliazione delle scorte rese o distrutte e verifica della ricezione degli avvisi presso i distributori, soprattutto quando la portata è multi-stato o internazionale.
    \item \textbf{Controlli integrati su catena del freddo e linee di confezionamento}: l'esposizione di \textit{Plastica / Busta} e \textit{Vetro / Bottiglia / Barattolo} a richiami per \textit{Salmonella}, \textit{Listeria} e allergeni non dichiarati richiama l'esigenza di rafforzare i punti critici di controllo (HACCP) proprio nelle fasi finali di riempimento, sigillatura ed etichettatura, oltre al monitoraggio continuo delle escursioni termiche nei segmenti \textit{Refrigerato / Fresco} e \textit{Surgelato}.
\end{enumerate}

\subsection{Limiti Metodologici}
Una valutazione rigorosa dei risultati richiede di esplicitare quattro limiti intrinseci al dataset e alla pipeline analitica adottata:
\begin{itemize}
    \item \textbf{Tassonomia basata su regole testuali e parole chiave}: i campi calcolati \texttt{Categoria Ragione di Richiamo}, \texttt{Tipo Packaging}, \texttt{Modalità Conservazione}, \texttt{Portata Distributiva} e \texttt{Dettaglio Allergene} operano tramite corrispondenze lessicali su stringhe in linguaggio naturale. Pur essendo efficaci per l'aggregazione esplorativa, tali regole possono risentire di sinonimi non mappati, descrizioni ambigue o compresenza di più concetti nello stesso testo.
    \item \textbf{Ampiezza della classe residuale ``Altro / Non specificato''}: nei campi derivati da testo libero (in particolare nella conservazione a temperatura ambiente/altro e nel packaging), una quota significativa di record non contiene indicazioni esplicite nel testo sorgente e confluisce nella categoria residuale. \todo{Quantificare esattamente l'incidenza percentuale della classe residuale ``Altro / Non specificato'' per Tipo Packaging, Modalità Conservazione e Categoria Ragione di Richiamo.}
    \item \textbf{Natura amministrativa e auto-dichiarata dei dati FDA}: il registro riflette tempistiche di classificazione (\texttt{Center Classification Date}) e aggiornamenti di stato (\texttt{Status}) che possono subire ritardi amministrativi rispetto al momento effettivo del ritiro fisico dal mercato. Inoltre, l'unità di conteggio adottata (\texttt{Conteggio prodotti}) pondera maggiormente eventi multi-referenza rispetto a richiami mono-prodotto.
    \item \textbf{Residua frammentazione anagrafica}: nonostante il campo \texttt{Ragione Sociale (Pulita)}, operatori multi-sito come \textit{Medline Industries LP} e \textit{Medline Industries LP - Northfield} possono ancora apparire come entità distinte nella classifica Top~15 quando i record mantengono riferimenti di stabilimento differenti.
\end{itemize}
